\documentclass[11pt, a4paper]{article}

\usepackage[utf8]{inputenc}
\usepackage[T1]{fontenc}
\usepackage[french]{babel}
\usepackage{amsmath, amssymb}
\usepackage{geometry}
\usepackage{graphicx}
\usepackage{tabularx}
\usepackage{booktabs}
\usepackage{xcolor}

% Configuration des marges
\geometry{
    top=2cm,
    bottom=2cm,
    left=1.5cm,
    right=1.5cm
}

% Commande pour les placeholders d'images
\newcommand{\imageplaceholder}[2]{
    \begin{center}
        \fbox{\begin{minipage}{#1}\centering\vspace{1cm}\textit{#2}\vspace{1cm}\end{minipage}}
    \end{center}
}

% Commande pour afficher la correction avec explication
\newcommand{\corr}[1]{
    \vspace{0.3cm}
    \begin{center}
    \fcolorbox{blue}{blue!5}{
    \begin{minipage}{0.95\textwidth}
    \color{blue} \textbf{Correction \& Raisonnement :} \\
    #1
    \end{minipage}}
    \end{center}
    \vspace{0.3cm}
}

\begin{document}

% ==================== EN-TÊTE ====================
\noindent
\begin{tabularx}{\textwidth}{@{}X c r@{}}
\textbf{EP361} & \textbf{vendredi 27 Juin 2025} & \textbf{Sans document} \\
\textbf{Electronique} & \textbf{Examen 1h45} & \textbf{Avec calculette} \\
\end{tabularx}
\vspace{0.2cm}
\hrule
\vspace{0.4cm}

% ==================== PROBLÈME I ====================
\begin{center}
    \textbf{\Large Problème I Filtre} \hfill \textbf{20 pts}
\end{center}
\vspace{-0.2cm}
\hrule
\vspace{0.3cm}

\subsubsection*{Partie A: Quadripôles - 10 points}

\imageplaceholder{0.5\textwidth}{[Insérer Figure 1 : Cellule en T]}
\begin{center}
    {\small Fig. 1 : Cellule en T.}
\end{center}

\noindent a) Exprimez $Y_{s}$ la matrice admittance d'un élément série $y_{o}$.

\corr{
\textit{Méthode :} Un dipôle d'admittance $y_o$ monté en série entre le port 1 et le port 2 d'un quadripôle obéit aux lois de Kirchhoff. La tension à ses bornes est $(v_1 - v_2)$.\\
Le courant entrant au port 1 est : $i_1 = y_o(v_1 - v_2) = y_o v_1 - y_o v_2$.\\
Le courant entrant au port 2 est : $i_2 = y_o(v_2 - v_1) = -y_o v_1 + y_o v_2$.\\
On peut l'écrire sous forme matricielle :
$$ \begin{bmatrix} i_1 \\ i_2 \end{bmatrix} = \begin{bmatrix} y_o & -y_o \\ -y_o & y_o \end{bmatrix} \begin{bmatrix} v_1 \\ v_2 \end{bmatrix} \quad \implies \quad Y_s = \begin{bmatrix} y_o & -y_o \\ -y_o & y_o \end{bmatrix} $$
}

\noindent b) Exprimez les coefficients $y_{ij}$ de la matrice admittance $Y_{T}$ en fonction des admittances $y_{1}$, $y_2$, $y_3$ des dipôles Fig. 1.

\corr{
\textit{Méthode (Théorème de Millman / Loi des nœuds) :} Soit $v_x$ la tension au nœud central du T.
La loi des nœuds en ce point donne : $i_1 + i_2 = y_3 v_x$, mais on peut aussi l'écrire avec les tensions extérieures.
Exprimons $v_x$ avec le théorème de Millman :
$$ v_x = \frac{y_1 v_1 + y_2 v_2}{y_1 + y_2 + y_3} $$
Les courants d'entrée et de sortie s'écrivent :
$$ i_1 = y_1(v_1 - v_x) \quad \text{et} \quad i_2 = y_2(v_2 - v_x) $$
En remplaçant $v_x$ dans l'expression de $i_1$ :
$$ i_1 = y_1 v_1 - y_1 \frac{y_1 v_1 + y_2 v_2}{y_1 + y_2 + y_3} = \frac{y_1(y_1 + y_2 + y_3) - y_1^2}{y_1 + y_2 + y_3} v_1 - \frac{y_1 y_2}{y_1 + y_2 + y_3} v_2 = \frac{y_1(y_2 + y_3)}{y_1 + y_2 + y_3} v_1 - \frac{y_1 y_2}{y_1 + y_2 + y_3} v_2 $$
Par symétrie, pour $i_2$ on obtient :
$$ i_2 = \frac{-y_1 y_2}{y_1 + y_2 + y_3} v_1 + \frac{y_2(y_1 + y_3)}{y_1 + y_2 + y_3} v_2 $$
La matrice admittance s'écrit donc :
$$ Y_T = \frac{1}{y_1 + y_2 + y_3} \begin{bmatrix} y_1(y_2 + y_3) & -y_1 y_2 \\ -y_1 y_2 & y_2(y_1 + y_3) \end{bmatrix} $$
}

\noindent c) Le quadripôle est chargé par un dipôle d'admittance $y_{u}$, exprimez $T(p)$ en fonction des $y_{ij}$ et de $y_{u}$ : $T(p) = \frac{v_2}{v_1}$.

\corr{
\textit{Méthode :} En connectant une charge $y_u$ au port 2, le courant $i_2$ s'exprime par $i_2 = -y_u v_2$ (le signe moins vient de la convention récepteur du quadripôle où $i_2$ est entrant).
D'après la définition de la matrice admittance :
$$ i_2 = y_{21} v_1 + y_{22} v_2 $$
On égalise les deux expressions :
$$ -y_u v_2 = y_{21} v_1 + y_{22} v_2 \implies v_2 (y_u + y_{22}) = -y_{21} v_1 $$
Ce qui donne la fonction de transfert :
$$ T(p) = \frac{v_2}{v_1} = \frac{-y_{21}}{y_{22} + y_u} $$
}

\vspace{0.5cm}
\subsubsection*{Partie B: Filtre à quadripôles - 10 pts}

Le filtre ci-dessous est constitué d'un AOP parfait, alimenté en symétrique associé à deux quadripôles $Q_{A}, Q_{B}$. Les quadripôles A et B sont décrits par leur matrice admittance.

\begin{minipage}{0.45\textwidth}
\imageplaceholder{0.9\textwidth}{[Insérer Figure 2 : Schéma AOP]}
\begin{center}{\small Fig. 2}\end{center}
\end{minipage}
\begin{minipage}{0.5\textwidth}
$$ Y_A = \begin{bmatrix} y_{A_{11}} & y_{A_{12}} \\ y_{A_{21}} & y_{A_{22}} \end{bmatrix} $$
$$ Y_B = \begin{bmatrix} y_{B_{11}} & y_{B_{12}} \\ y_{B_{21}} & y_{B_{22}} \end{bmatrix} $$
\end{minipage}

\textbf{Données :}\\
Quadripôle A - Matrice Admittance : $Y_A = \frac{1}{y_1 + y_2 + y_3} \begin{bmatrix} y_1(y_2 + y_3) & -y_1 y_2 \\ -y_1 y_2 & y_2(y_1 + y_3) \end{bmatrix}$ \\
Quadripôle B - Matrice Admittance : $Y_B = \begin{bmatrix} \frac{y_1(y_2+y_3)}{y_1+y_2+y_3} + y_4 & \frac{-y_1 y_2}{y_1+y_2+y_3} - y_4 \\[10pt] \frac{-y_1 y_2}{y_1+y_2+y_3} - y_4 & \frac{y_2(y_1+y_3)}{y_1+y_2+y_3} + y_4 \end{bmatrix}$

\vspace{0.3cm}
\noindent a/ Exprimez en démontrant $H(p)$ en fonction de certains éléments $y_{x_{ij}} : H(p)=\frac{v_s}{v_e}$

\corr{
\textit{Méthode (AOP en régime linéaire + Masse virtuelle) :}
L'AOP est supposé parfait et en boucle fermée sur l'entrée inverseuse ($V^-$). La tension différentielle est donc nulle : $V^+ = V^- = 0$. De plus, le courant d'entrée de l'AOP est nul ($i^- = 0$).
\begin{itemize}
    \item La sortie de $Q_A$ est reliée à $V^-$. Le courant sortant de $Q_A$ est $i_{A2} = y_{A21}v_e + y_{A22}V^-$. Puisque $V^-=0$, on a $i_{A2} = y_{A21}v_e$.
    \item L'entrée de $Q_B$ est reliée à $V^-$ et sa sortie à $v_s$. Le courant entrant dans $Q_B$ est $i_{B1} = y_{B11}V^- + y_{B12}v_s$. Puisque $V^-=0$, on a $i_{B1} = y_{B12}v_s$.
\end{itemize}
Loi des nœuds en $V^-$ : la somme des courants est nulle $\implies i_{A2} + i_{B1} = 0$.
$$ y_{A21} v_e + y_{B12} v_s = 0 \implies H(p) = \frac{v_s}{v_e} = -\frac{y_{A21}}{y_{B12}} $$
}

\noindent b/ Exprimez $H(p)$ la fonction de transfert du filtre Fig. 3 :\\
$y_1$ : résistance R \qquad $y_3$ : capacité C\\
$y_2$ : résistance R \qquad $y_4$ : capacité C

\corr{
D'après les matrices données :
$$ y_{A21} = \frac{-y_1 y_2}{y_1 + y_2 + y_3} \quad \text{et} \quad y_{B12} = \frac{-y_1 y_2}{y_1 + y_2 + y_3} - y_4 $$
On utilise l'expression de $H(p)$ :
$$ H(p) = -\frac{\frac{-y_1 y_2}{y_1 + y_2 + y_3}}{\frac{-y_1 y_2}{y_1 + y_2 + y_3} - y_4} = \frac{y_1 y_2}{-y_1 y_2 - y_4(y_1 + y_2 + y_3)} $$
On remplace par les admittances physiques ($y_1 = y_2 = 1/R$ et $y_3 = y_4 = Cp$) :
$$ H(p) = \frac{(1/R)(1/R)}{-(1/R)(1/R) - Cp(1/R + 1/R + Cp)} = \frac{1/R^2}{-1/R^2 - 2Cp/R - C^2 p^2} $$
On multiplie en haut et en bas par $-R^2$ pour normaliser :
$$ H(p) = -\frac{1}{1 + 2RCp + (RCp)^2} = -\frac{1}{(1 + RCp)^2} $$
}

\noindent c/ Définissez l'ordre et le type de filtre.

\corr{
\textit{Analyse de la fonction de transfert :}
La fonction de transfert est $H(p) = -\frac{1}{(1 + RCp)^2}$.
\begin{itemize}
    \item \textbf{Type :} C'est un filtre \textbf{passe-bas}. (Le gain vaut 1 en basse fréquence $p=0$ et tend vers 0 en haute fréquence $p \rightarrow \infty$).
    \item \textbf{Ordre :} Le dénominateur est un polynôme de degré 2 en $p$. C'est un filtre du \textbf{2ème ordre} (plus précisément, deux pôles réels identiques).
\end{itemize}
}

\vspace{0.5cm}
% ==================== PROBLÈME II ====================
\begin{center}
    \textbf{\Large Problème II Filtre} \hfill \textbf{10 pts}
\end{center}
\vspace{-0.2cm}
\hrule
\vspace{0.3cm}

\imageplaceholder{0.7\textwidth}{[Insérer Figure 3 : Filtre AOP multiples]}

\noindent Les AOP sont parfaits alimentés symétriques $+/-V_{cc}$.
\begin{itemize}
    \item Exprimez la tension $v_2$ en fonction de $v_e$ et $v_1$.
    \item Exprimez la tension $v_s$ en fonction de $v_2$.
    \item Exprimez la tension $v_1$ en fonction de $v_s$.
    \item Exprimez $T(p)$ la fonction de transfert sous forme canonique en exprimant la pulsation propre, le coefficient d'amortissement et le gain statique : $T(p) = \frac{v_s}{v_e}$.
    \item Quel est le type et l'ordre du filtre ?
    \item Pourquoi tant de composants pour ce filtre ?
\end{itemize}

\corr{
\textit{Analyse des étages du filtre universel (variables d'état) :}\\
1. \textbf{AOP $A_1$ (Sommateur-Intégrateur ou Dérivateur) :}
L'entrée $v_e$ attaque une capacité C, et $v_1$ attaque une autre capacité C. La contre-réaction est une résistance R. Le théorème de Millman à l'entrée inverseuse de $A_1$ ($V^- = 0$) donne :
$$ \frac{v_e \cdot Cp + v_1 \cdot Cp + v_2 / R}{Cp + Cp + 1/R} = 0 \implies v_2 = -RCp(v_e + v_1) $$

2. \textbf{AOP $A_3$ (Inverseur unitaire) :}
C'est un simple amplificateur inverseur avec deux résistances $R_1$ identiques :
$$ v_1 = -\frac{R_1}{R_1} v_s \implies v_1 = -v_s $$

3. \textbf{AOP $A_2$ (Intégrateur à fuite / Filtre passe-bas) :}
L'entrée $v_2$ attaque une capacité C. La boucle de rétroaction contient R en parallèle avec $C_3$. L'admittance de retour est $Y_{ret} = 1/R + C_3 p$.
$$ v_s = - \frac{Y_{in}}{Y_{ret}} v_2 = - \frac{Cp}{1/R + C_3 p} v_2 = - \frac{RCp}{1 + R C_3 p} v_2 $$

4. \textbf{Fonction de transfert globale $T(p)$ :}
On remplace $v_1$ par $-v_s$ dans l'équation 1 : $v_2 = -RCp(v_e - v_s)$.
On insère cela dans l'équation 3 :
$$ v_s = - \left( -RCp(v_e - v_s) \right) \frac{RCp}{1 + R C_3 p} = (v_e - v_s) \frac{(RCp)^2}{1 + R C_3 p} $$
On regroupe les termes en $v_s$ :
$$ v_s \left[ 1 + \frac{(RCp)^2}{1 + R C_3 p} \right] = v_e \frac{(RCp)^2}{1 + R C_3 p} $$
$$ T(p) = \frac{v_s}{v_e} = \frac{(RCp)^2}{1 + R C_3 p + (RCp)^2} $$

5. \textbf{Identification canonique, Type et Ordre :}
C'est une structure du type $H(p) = \frac{A_0 \cdot (p/\omega_0)^2}{1 + 2m(p/\omega_0) + (p/\omega_0)^2}$.
\begin{itemize}
    \item \textbf{Type :} \textbf{Passe-haut} (le numérateur est en $p^2$).
    \item \textbf{Ordre :} \textbf{2ème ordre}.
    \item Pulsation propre : $\omega_0 = \frac{1}{RC}$.
    \item Gain statique (haute fréquence ici) : $A_0 = 1$.
    \item Coefficient d'amortissement : Le terme du 1er degré est $RC_3 p$. On identifie $\frac{2m}{\omega_0} = RC_3 \implies 2m RC = RC_3 \implies m = \frac{C_3}{2C}$.
\end{itemize}

6. \textbf{Pourquoi tant de composants ?}
Ce type de montage (filtre à variables d'état) permet de régler indépendamment la fréquence de coupure $\omega_0$ (en modifiant R) et le facteur d'amortissement $m$ (en modifiant $C_3$). Cela offre une souplesse de réglage impossible avec une simple cellule de Sallen-Key.
}

\vspace{0.5cm}
% ==================== PROBLÈME III ====================
\begin{center}
    \textbf{\Large Problème III AOP de puissance} \hfill \textbf{10 pts}
\end{center}
\vspace{-0.2cm}
\hrule
\vspace{0.3cm}

\imageplaceholder{0.7\textwidth}{[Insérer Figure 4 : Ampli de Puissance]}

La tension $v_1(t)$ est un signal triangulaire symétrique, imposé par l'extérieur, d'amplitude crête $\hat{v}_1$ de 0,3 V de fréquence 1 kHz, centré en zéro.
\begin{itemize}
    \item AOP parfait. Transistors $T_1$ et $T_2$ : $\beta = 100 \gg 1$, $|v_{be}| = 0.7 V$.
    \item Résistances : $R_1 = 1 k\Omega$, $R_2 = 9 k\Omega$, $R_u = 8 \Omega$. Alim $\pm V_{cc} = \pm 15 V$.
\end{itemize}

\noindent Exprimez l'amplification en tension $A=v_2/v_1$. Calculez A. \\
Calculez $\hat{v}_2$ la tension crête aux bornes de $R_u$. \\
Calculez les valeurs crêtes de $I_u$ et $I_{out}$. \\
Tracé $v_2(t)$, $V_{BM}(t)$ sur la grille prévue. \\
Exprimez la puissance active (moyenne) $P_u$ dans $R_u$ en fonction de $\hat{v}_2$. Calculez $P_u$.

\corr{
\textbf{1. Amplification en tension et tension de sortie crête :}\\
L'AOP est monté en amplificateur non-inverseur. L'étage de puissance push-pull (T1/T2) est inséré \textit{dans} la boucle de contre-réaction. Par conséquent, l'AOP compense les chutes de tension $v_{be}$.
Le diviseur de tension sur le retour fixe $V^- = v_2 \frac{R_1}{R_1 + R_2}$. Comme l'AOP est parfait, $V^+ = V^- \implies v_1 = v_2 \frac{R_1}{R_1 + R_2}$.
$$ A = \frac{v_2}{v_1} = \frac{R_1 + R_2}{R_1} = 1 + \frac{9k}{1k} = \mathbf{10} $$
La tension de sortie crête est $\hat{v}_2 = A \times \hat{v}_1 = 10 \times 0.3V = \mathbf{3 V}$.

\textbf{2. Courants crêtes $I_u$ et $I_{out}$ :}\\
Le courant maximum dans la charge $R_u$ est $\hat{I}_u = \frac{\hat{v}_2}{R_u} = \frac{3V}{8\Omega} = \mathbf{0.375 A}$ ($375 mA$).
Le courant $i_{out}$ sortant de l'AOP doit fournir le courant de base du transistor et le courant de la boucle de retour.
Courant de retour : $\hat{I}_{ret} = \frac{\hat{v}_2}{R_1 + R_2} = \frac{3V}{10k\Omega} = 0.3 mA$.
Courant de base : Le transistor T1 fournit $\hat{I}_u$. Donc $\hat{i}_B \approx \frac{\hat{I}_u}{\beta} = \frac{375mA}{100} = 3.75 mA$.
Le courant crête total fourni par l'AOP est $\hat{I}_{out} = 3.75 + 0.3 = \mathbf{4.05 mA}$. (Cela reste largement dans les capacités d'un AOP classique).

\textbf{3. Tracé de $v_2(t)$ et $V_{BM}(t)$ :}\\
$v_2(t)$ est un signal triangulaire parfaitement proportionnel à $v_1(t)$ avec une amplitude de $\pm 3V$.\\
$V_{BM}(t)$ est la tension à la sortie de l'AOP (base des transistors). 
Lorsque $v_2 > 0$, T1 conduit, il faut que $V_{BM} = v_2 + v_{BE1} = v_2 + 0.7V$. L'AOP "saute" instantanément à $+0.7V$ lorsque le signal croise zéro pour éliminer la distorsion de croisement.
Lorsque $v_2 < 0$, T2 conduit, $V_{BM} = v_2 - 0.7V$.
(\textit{Sur la grille, $V_{BM}$ suit la forme du triangle mais décalé de 0.7V vers le haut pour les phases positives et de 0.7V vers le bas pour les phases négatives, créant un saut vertical au passage par zéro}).

\textbf{4. Puissance active $P_u$ :}\\
Pour un signal triangulaire d'amplitude crête $\hat{v}_2$, la valeur efficace est $V_{eff} = \frac{\hat{v}_2}{\sqrt{3}}$.
La puissance moyenne dissipée dans une résistance est :
$$ P_u = \frac{V_{eff}^2}{R_u} = \frac{\left( \hat{v}_2 / \sqrt{3} \right)^2}{R_u} = \frac{\hat{v}_2^2}{3 R_u} $$
Calcul numérique :
$$ P_u = \frac{3^2}{3 \times 8} = \frac{9}{24} = \frac{3}{8} = \mathbf{0.375 W} \quad \text{(soit 375 mW)} $$
}

\vfill
% ==================== PIED DE PAGE ====================
\hrule
\vspace{0.1cm}
\noindent
Dehay \hfill Esisar \hfill 1/2

\end{document}